% Smaller XOR circuits for AES MixColumns -- short note for the IACR ePrint Archive.
% Build: python3 generate_appendix.py && tectonic mixcolumns_note.tex
\documentclass[11pt,a4paper]{article}
\usepackage[margin=2.7cm]{geometry}
\usepackage{amsmath,amssymb}
\usepackage{booktabs}
\usepackage{multicol}
\usepackage[hidelinks]{hyperref}
\urlstyle{tt}

\title{Smaller XOR Circuits for AES MixColumns:\\
89 Gates; 98 Gates at Depth 3; 91 Gates at Depth 6}
\author{Joe Bachir\\[2pt]
Independent researcher\\[2pt]
\texttt{joebachir20@gmail.com}}
\date{July 2026}

\begin{document}
\maketitle

\begin{abstract}
We report three explicit implementations of the AES MixColumns linear
transformation as circuits of 2-input XOR gates over $\mathrm{GF}(2)$: an
\textbf{89-gate} circuit of depth 10; a \textbf{98-gate} circuit of depth 3;
and a \textbf{91-gate} circuit of depth 6. The smallest previously published
XOR count we are aware of is 91, by Lin, Xiang, Zeng, and Zhang (CT-RSA 2021),
stated in the s-XOR model, whose programs translate instruction-for-gate into
2-input XOR circuits; the classic baseline stated directly in the 2-input-XOR
model is the 92-gate, depth-6 circuit of Maximov (ePrint 2019/833). The
89-gate circuit improves on both. At depth 3 --- the known minimum depth for
this map --- the 98-gate circuit improves on the 99-gate result of Shi, Feng,
and Xu (ToSC 2023). The 91-gate circuit matches the smallest published count
while achieving depth 6. All three circuits are given in full in the
appendices and are published as machine-checkable artifacts with self-contained
verifiers that rebuild the MixColumns specification from scratch. We make no
optimality claim on gate counts.
\end{abstract}

\section{Model and specification}

A circuit is an ordered list of 2-input XOR gates. Signals $s_0,\dots,s_{31}$
are the input bits; gate $k$ produces signal $s_{32+k} = s_a \oplus s_b$ with
$a,b$ both strictly smaller than $32+k$. Depth counts gates on the longest
input-to-signal path, with inputs at depth 0.

The bit/byte convention is: the AES state column is four bytes
$a_0,a_1,a_2,a_3$; bit index $i \in \{0,\dots,31\}$ denotes bit $(i \bmod 8)$
of byte $\lfloor i/8 \rfloor$, least-significant-bit first within each byte.
MixColumns maps the input column to the output column with
\[
\mathit{out}[c] \;=\; 2 \cdot a_c \,\oplus\, 3 \cdot a_{(c+1) \bmod 4}
\,\oplus\, a_{(c+2) \bmod 4} \,\oplus\, a_{(c+3) \bmod 4},
\]
with multiplication in $\mathrm{GF}(2^8)$ modulo $x^8+x^4+x^3+x+1$
(\texttt{0x11b}), following NIST FIPS 197-upd1, Section 5.1.3,
Eq.~5.6~\cite{FIPS197}. Because MixColumns is linear over $\mathrm{GF}(2)$,
agreement on the 32 unit-input vectors is a complete correctness check. Gate
counts and depths are invariant under input/output bit relabeling, so the
comparisons below do not depend on convention choices.

\section{Results}

\begin{table}[h]
\centering
\begin{tabular}{lccl}
\toprule
Circuit & Gates & Depth & Published comparison figure \\
\midrule
\texttt{mixcolumns\_89gates}        & \textbf{89} & 10 & 91~\cite{LXZZ21} (s-XOR); 92~\cite{Max19} \\
\texttt{mixcolumns\_98gates\_depth3} & \textbf{98} & 3  & 99 at depth 3~\cite{SFX23} \\
\texttt{mixcolumns\_91gates\_depth6} & 91 & \textbf{6} & 92 at depth 6~\cite{Max19}; ties 91~\cite{LXZZ21} \\
\bottomrule
\end{tabular}
\caption{The three circuits (full listings in Appendices~\ref{app:mixcolumns_89gates}--\ref{app:mixcolumns_91gates_depth6}).}
\end{table}

The smallest previously published XOR count we are aware of is 91, by Lin,
Xiang, Zeng, and Zhang~\cite{LXZZ21}, stated in the s-XOR (in-place) model.
Every $k$-instruction s-XOR program yields a $k$-gate 2-input XOR circuit
(each in-place instruction $x_i \leftarrow x_i \oplus x_j$ is one 2-input XOR
gate), so we compare against 91 rather than only against Maximov's explicit
92-gate circuit~\cite{Max19}. Yuan et al.~\cite{YWSZZ24} likewise report 91
XORs in an s-XOR / quantum-depth framing.

Two remarks. First, depth 3 is the known minimum depth for AES MixColumns
(stated e.g.\ in~\cite{SFX23}; an output depending on $w$ inputs needs depth
at least $\lceil \log_2 w \rceil$, and MixColumns has outputs of weight~7).
The contribution of the depth-3 circuit is therefore the gate count at that
depth, not the depth itself. Second, the 91-gate depth-6 circuit matches the
smallest published gate count while using one gate fewer than Maximov's
depth-6 circuit; we include it as an explicit shallow trade-off point between
the 89-gate circuit and the depth-3 circuit.

\section{Verification}
\label{sec:verification}

The circuits are static, machine-checkable artifacts, published with
verifiers at
\begin{center}
\url{https://github.com/Joe-b-20/aes-mixcolumns-xor-circuits}
\end{center}
and archived permanently at DOI
\href{https://doi.org/10.5281/zenodo.21299092}{10.5281/zenodo.21299092}.
Each circuit is verified by two separately implemented, dependency-free Python
paths and one hardware path:

\begin{enumerate}
\item \texttt{verify.py} rebuilds the $32 \times 32$ $\mathrm{GF}(2)$ target
  map from the $\mathrm{GF}(2^8)$ definition above --- the verifier
  \emph{is} the specification, so there is no hidden convention to mismatch
  --- and checks structure (every gate a 2-input XOR of earlier signals),
  all 32 output masks, the declared gate count and depth, and the SHA-256
  integrity metadata.
\item \texttt{audit/cleanroom\_verify.py}, a separately written clean-room
  verifier, rechecks everything against an independent byte-level reference,
  runs $10^5$ deterministic random trials per circuit, verifies there are no
  dead gates or duplicate intermediate values, and demonstrates that mutated
  or convention-mismatched artifacts are rejected (14 adversarial cases).
\item Verilog netlists and testbenches for all three circuits pass under
  Icarus Verilog, checking all 32 basis vectors in simulation; a continuous
  integration workflow runs every path on each commit.
\end{enumerate}

For integrity, each appendix states the circuit's canonical SHA-256, computed
over the UTF-8 bytes of the compact JSON string
\texttt{\{"inputCount":32,"gates":[...]\}}; the repository's
\texttt{bounds.json} carries the same values, and both verifiers recheck them.

\section{Scope and claims}

\begin{itemize}
\item \textbf{Not optimality.} Minimum 2-input-XOR circuit size (the Shortest
  Linear Program problem) is NP-hard; we do not prove any of these counts
  minimal. We claim only that they are the smallest we have found or seen
  published as of July 2026.
\item \textbf{Comparisons.} All counts are for the exact model and convention
  of Section~1; the baselines~\cite{Max19,LXZZ21,SFX23,YWSZZ24} were
  source-checked. Results in other cost models (multi-input XOR gates,
  gate-equivalent area, technology-mapped netlists, in-place quantum CNOT
  circuits) are not comparable and are not claimed against.
\item \textbf{Method.} The search technique used to find these circuits is
  deferred to separate, ongoing work; none of the claims here depend on it.
  The artifacts are self-contained and permanently verifiable.
\end{itemize}

\begin{thebibliography}{9}

\bibitem{FIPS197}
National Institute of Standards and Technology,
\emph{Advanced Encryption Standard (AES)},
NIST FIPS 197-upd1, May 2023.
\url{https://doi.org/10.6028/NIST.FIPS.197-upd1}

\bibitem{Max19}
A.~Maximov,
\emph{AES MixColumn with 92 XOR gates},
IACR Cryptology ePrint Archive, Report 2019/833, 2019.
\url{https://eprint.iacr.org/2019/833}

\bibitem{LXZZ21}
D.~Lin, Z.~Xiang, X.~Zeng, and S.~Zhang,
\emph{A Framework to Optimize Implementations of Matrices},
Topics in Cryptology -- CT-RSA 2021, LNCS 12704, Springer, 2021.
\url{https://doi.org/10.1007/978-3-030-75539-3_25}

\bibitem{SFX23}
H.~Shi, X.~Feng, and S.~Xu,
\emph{A Framework with Improved Heuristics to Optimize Low-Latency
Implementations of Linear Layers},
IACR Transactions on Symmetric Cryptology, 2023(4):489--510, 2023.
\url{https://doi.org/10.46586/tosc.v2023.i4.489-510}

\bibitem{YWSZZ24}
Y.~Yuan, W.~Wu, T.~Shi, L.~Zhang, and Y.~Zhang,
\emph{A Framework to Improve the Implementations of Linear Layers},
IACR Transactions on Symmetric Cryptology, 2024(2):322--347, 2024.
\url{https://doi.org/10.46586/tosc.v2024.i2.322-347}

\end{thebibliography}

\appendix
% Auto-generated by paper/generate_appendix.py -- do not edit by hand.
\section{Circuit \texttt{mixcolumns\_97gates\_depth3} (97 gates, depth 3)}
\label{app:mixcolumns_97gates_depth3}
Signals $s_0,\dots,s_{31}$ are the inputs; gate $k$ produces $s_{32+k}$. Canonical gate-list SHA-256 (see Section~\ref{sec:verification}): \\
\texttt{\small e4c4e678c4b7c22d889c6e060b43859c111b82b953d6d9b43c49c079c0545b56}
\begin{multicols}{3}\scriptsize\ttfamily
s32 = s0 $\oplus$ s8 \\
s33 = s5 $\oplus$ s13 \\
s34 = s6 $\oplus$ s14 \\
s35 = s7 $\oplus$ s15 \\
s36 = s8 $\oplus$ s16 \\
s37 = s9 $\oplus$ s17 \\
s38 = s10 $\oplus$ s18 \\
s39 = s11 $\oplus$ s19 \\
s40 = s11 $\oplus$ s20 \\
s41 = s12 $\oplus$ s20 \\
s42 = s6 $\oplus$ s22 \\
s43 = s15 $\oplus$ s23 \\
s44 = s1 $\oplus$ s24 \\
s45 = s1 $\oplus$ s25 \\
s46 = s9 $\oplus$ s25 \\
s47 = s2 $\oplus$ s26 \\
s48 = s11 $\oplus$ s26 \\
s49 = s2 $\oplus$ s27 \\
s50 = s3 $\oplus$ s27 \\
s51 = s4 $\oplus$ s27 \\
s52 = s4 $\oplus$ s28 \\
s53 = s12 $\oplus$ s28 \\
s54 = s13 $\oplus$ s29 \\
s55 = s21 $\oplus$ s29 \\
s56 = s22 $\oplus$ s30 \\
s57 = s7 $\oplus$ s31 \\
s58 = s23 $\oplus$ s31 \\
s59 = s33 $\oplus$ s14 \\
s60 = s35 $\oplus$ s17 \\
s61 = s39 $\oplus$ s10 \\
s62 = s41 $\oplus$ s5 \\
s63 = s35 $\oplus$ s40 \\
s64 = s42 $\oplus$ s7 \\
s65 = s36 $\oplus$ s43 \\
s66 = s43 $\oplus$ s19 \\
s67 = s32 $\oplus$ s24 \\
s68 = s36 $\oplus$ s24 \\
s69 = s44 $\oplus$ s37 \\
s70 = s32 $\oplus$ s46 \\
s71 = s46 $\oplus$ s10 \\
s72 = s46 $\oplus$ s16 \\
s73 = s45 $\oplus$ s17 \\
s74 = s45 $\oplus$ s38 \\
s75 = s47 $\oplus$ s37 \\
s76 = s47 $\oplus$ s39 \\
s77 = s46 $\oplus$ s26 \\
s78 = s49 $\oplus$ s35 \\
s79 = s50 $\oplus$ s18 \\
s80 = s50 $\oplus$ s38 \\
s81 = s51 $\oplus$ s19 \\
s82 = s51 $\oplus$ s41 \\
s83 = s53 $\oplus$ s3 \\
s84 = s52 $\oplus$ s40 \\
s85 = s52 $\oplus$ s21 \\
s86 = s42 $\oplus$ s54 \\
s87 = s53 $\oplus$ s54 \\
s88 = s43 $\oplus$ s56 \\
s89 = s55 $\oplus$ s30 \\
s90 = s57 $\oplus$ s0 \\
s91 = s57 $\oplus$ s3 \\
s92 = s57 $\oplus$ s15 \\
s93 = s34 $\oplus$ s58 \\
s94 = s44 $\oplus$ s58 \\
s95 = s48 $\oplus$ s58 \\
s96 = s53 $\oplus$ s58 \\
s97 = s68 $\oplus$ s35 \\
s98 = s70 $\oplus$ s60 \\
s99 = s74 $\oplus$ s77 \\
s100 = s78 $\oplus$ s61 \\
s101 = s63 $\oplus$ s83 \\
s102 = s87 $\oplus$ s85 \\
s103 = s59 $\oplus$ s56 \\
s104 = s93 $\oplus$ s15 \\
s105 = s65 $\oplus$ s67 \\
s106 = s65 $\oplus$ s73 \\
s107 = s75 $\oplus$ s18 \\
s108 = s80 $\oplus$ s66 \\
s109 = s84 $\oplus$ s66 \\
s110 = s62 $\oplus$ s55 \\
s111 = s86 $\oplus$ s89 \\
s112 = s93 $\oplus$ s64 \\
s113 = s67 $\oplus$ s58 \\
s114 = s94 $\oplus$ s72 \\
s115 = s75 $\oplus$ s71 \\
s116 = s95 $\oplus$ s79 \\
s117 = s96 $\oplus$ s81 \\
s118 = s87 $\oplus$ s62 \\
s119 = s89 $\oplus$ s34 \\
s120 = s92 $\oplus$ s56 \\
s121 = s90 $\oplus$ s36 \\
s122 = s69 $\oplus$ s90 \\
s123 = s74 $\oplus$ s2 \\
s124 = s76 $\oplus$ s91 \\
s125 = s82 $\oplus$ s91 \\
s126 = s85 $\oplus$ s33 \\
s127 = s86 $\oplus$ s59 \\
s128 = s88 $\oplus$ s64 \\
\end{multicols}
\noindent Output signals (output bit $j$ is carried by the listed signal, $j = 0,\dots,31$):
\begin{center}\scriptsize\ttfamily
[97, 98, 99, 100, 101, 102, 103, 104, 105, 106, 107, 108, 109, 110, 111, 112, 113, 114, 115, 116, 117, 118, 119, 120, 121, 122, 123, 124, 125, 126, 127, 128]
\end{center}

\section{Circuit \texttt{mixcolumns\_92gates\_depth4} (92 gates, depth 4)}
\label{app:mixcolumns_92gates_depth4}
Signals $s_0,\dots,s_{31}$ are the inputs; gate $k$ produces $s_{32+k}$. Canonical gate-list SHA-256 (see Section~\ref{sec:verification}): \\
\texttt{\small e79bfccf5322fe8933242b47b53e94e73163615c70d5c45c9721e27ac86fd97d}
\begin{multicols}{3}\scriptsize\ttfamily
s32 = s10 $\oplus$ s26 \\
s33 = s12 $\oplus$ s20 \\
s34 = s14 $\oplus$ s30 \\
s35 = s6 $\oplus$ s14 \\
s36 = s11 $\oplus$ s28 \\
s37 = s8 $\oplus$ s16 \\
s38 = s7 $\oplus$ s31 \\
s39 = s17 $\oplus$ s25 \\
s40 = s13 $\oplus$ s21 \\
s41 = s2 $\oplus$ s26 \\
s42 = s5 $\oplus$ s29 \\
s43 = s16 $\oplus$ s24 \\
s44 = s9 $\oplus$ s25 \\
s45 = s18 $\oplus$ s26 \\
s46 = s0 $\oplus$ s17 \\
s47 = s1 $\oplus$ s9 \\
s48 = s23 $\oplus$ s31 \\
s49 = s3 $\oplus$ s27 \\
s50 = s4 $\oplus$ s20 \\
s51 = s12 $\oplus$ s27 \\
s52 = s22 $\oplus$ s30 \\
s53 = s13 $\oplus$ s29 \\
s54 = s4 $\oplus$ s28 \\
s55 = s11 $\oplus$ s19 \\
s56 = s15 $\oplus$ s31 \\
s57 = s43 $\oplus$ s44 \\
s58 = s34 $\oplus$ s53 \\
s59 = s5 $\oplus$ s54 \\
s60 = s0 $\oplus$ s37 \\
s61 = s38 $\oplus$ s52 \\
s62 = s6 $\oplus$ s40 \\
s63 = s32 $\oplus$ s56 \\
s64 = s36 $\oplus$ s56 \\
s65 = s37 $\oplus$ s56 \\
s66 = s1 $\oplus$ s48 \\
s67 = s50 $\oplus$ s53 \\
s68 = s35 $\oplus$ s56 \\
s69 = s3 $\oplus$ s38 \\
s70 = s51 $\oplus$ s54 \\
s71 = s19 $\oplus$ s48 \\
s72 = s24 $\oplus$ s38 \\
s73 = s39 $\oplus$ s41 \\
s74 = s41 $\oplus$ s55 \\
s75 = s34 $\oplus$ s38 \\
s76 = s43 $\oplus$ s48 \\
s77 = s22 $\oplus$ s42 \\
s78 = s41 $\oplus$ s44 \\
s79 = s10 $\oplus$ s47 \\
s80 = s45 $\oplus$ s49 \\
s81 = s21 $\oplus$ s33 \\
s82 = s67 $\oplus$ s81 \\
s83 = s57 $\oplus$ s66 \\
s84 = s58 $\oplus$ s77 \\
s85 = s42 $\oplus$ s81 \\
s86 = s59 $\oplus$ s67 \\
s87 = s59 $\oplus$ s40 \\
s88 = s10 $\oplus$ s73 \\
s89 = s45 $\oplus$ s79 \\
s90 = s23 $\oplus$ s68 \\
s91 = s62 $\oplus$ s52 \\
s92 = s58 $\oplus$ s62 \\
s93 = s60 $\oplus$ s76 \\
s94 = s35 $\oplus$ s77 \\
s95 = s15 $\oplus$ s61 \\
s96 = s69 $\oplus$ s74 \\
s97 = s65 $\oplus$ s72 \\
s98 = s38 $\oplus$ s60 \\
s99 = s70 $\oplus$ s71 \\
s100 = s57 $\oplus$ s46 \\
s101 = s39 $\oplus$ s78 \\
s102 = s39 $\oplus$ s65 \\
s103 = s69 $\oplus$ s50 \\
s104 = s63 $\oplus$ s49 \\
s105 = s64 $\oplus$ s50 \\
s106 = s64 $\oplus$ s69 \\
s107 = s23 $\oplus$ s75 \\
s108 = s71 $\oplus$ s80 \\
s109 = s46 $\oplus$ s72 \\
s110 = s102 $\oplus$ s66 \\
s111 = s18 $\oplus$ s101 \\
s112 = s89 $\oplus$ s78 \\
s113 = s63 $\oplus$ s108 \\
s114 = s108 $\oplus$ s55 \\
s115 = s93 $\oplus$ s65 \\
s116 = s47 $\oplus$ s109 \\
s117 = s52 $\oplus$ s107 \\
s118 = s96 $\oplus$ s104 \\
s119 = s33 $\oplus$ s106 \\
s120 = s71 $\oplus$ s105 \\
s121 = s103 $\oplus$ s51 \\
s122 = s97 $\oplus$ s100 \\
s123 = s90 $\oplus$ s75 \\
\end{multicols}
\noindent Output signals (output bit $j$ is carried by the listed signal, $j = 0,\dots,31$):
\begin{center}\scriptsize\ttfamily
[97, 122, 89, 118, 119, 82, 84, 90, 115, 110, 111, 113, 120, 85, 91, 117, 93, 83, 88, 114, 99, 86, 92, 95, 98, 116, 112, 96, 121, 87, 94, 123]
\end{center}

\section{Circuit \texttt{mixcolumns\_89gates\_depth5} (89 gates, depth 5)}
\label{app:mixcolumns_89gates_depth5}
Signals $s_0,\dots,s_{31}$ are the inputs; gate $k$ produces $s_{32+k}$. Canonical gate-list SHA-256 (see Section~\ref{sec:verification}): \\
\texttt{\small c3b972af282d25e8313a510f5cf840472401c6268bd28cda889593fc3379f064}
\begin{multicols}{3}\scriptsize\ttfamily
s32 = s9 $\oplus$ s17 \\
s33 = s19 $\oplus$ s27 \\
s34 = s1 $\oplus$ s25 \\
s35 = s21 $\oplus$ s29 \\
s36 = s12 $\oplus$ s28 \\
s37 = s22 $\oplus$ s30 \\
s38 = s3 $\oplus$ s27 \\
s39 = s4 $\oplus$ s28 \\
s40 = s9 $\oplus$ s25 \\
s41 = s13 $\oplus$ s29 \\
s42 = s5 $\oplus$ s13 \\
s43 = s15 $\oplus$ s23 \\
s44 = s6 $\oplus$ s22 \\
s45 = s7 $\oplus$ s23 \\
s46 = s12 $\oplus$ s20 \\
s47 = s10 $\oplus$ s18 \\
s48 = s2 $\oplus$ s26 \\
s49 = s6 $\oplus$ s14 \\
s50 = s19 $\oplus$ s20 \\
s51 = s8 $\oplus$ s16 \\
s52 = s7 $\oplus$ s31 \\
s53 = s2 $\oplus$ s18 \\
s54 = s30 $\oplus$ s35 \\
s55 = s21 $\oplus$ s39 \\
s56 = s5 $\oplus$ s46 \\
s57 = s14 $\oplus$ s42 \\
s58 = s15 $\oplus$ s37 \\
s59 = s11 $\oplus$ s33 \\
s60 = s11 $\oplus$ s43 \\
s61 = s0 $\oplus$ s52 \\
s62 = s33 $\oplus$ s45 \\
s63 = s38 $\oplus$ s52 \\
s64 = s43 $\oplus$ s51 \\
s65 = s8 $\oplus$ s45 \\
s66 = s36 $\oplus$ s41 \\
s67 = s41 $\oplus$ s44 \\
s68 = s34 $\oplus$ s53 \\
s69 = s24 $\oplus$ s45 \\
s70 = s40 $\oplus$ s47 \\
s71 = s31 $\oplus$ s49 \\
s72 = s26 $\oplus$ s34 \\
s73 = s32 $\oplus$ s48 \\
s74 = s44 $\oplus$ s45 \\
s75 = s40 $\oplus$ s51 \\
s76 = s18 $\oplus$ s73 \\
s77 = s10 $\oplus$ s68 \\
s78 = s55 $\oplus$ s66 \\
s79 = s37 $\oplus$ s57 \\
s80 = s35 $\oplus$ s56 \\
s81 = s56 $\oplus$ s66 \\
s82 = s55 $\oplus$ s42 \\
s83 = s70 $\oplus$ s72 \\
s84 = s43 $\oplus$ s71 \\
s85 = s54 $\oplus$ s67 \\
s86 = s54 $\oplus$ s49 \\
s87 = s17 $\oplus$ s64 \\
s88 = s57 $\oplus$ s67 \\
s89 = s60 $\oplus$ s47 \\
s90 = s62 $\oplus$ s53 \\
s91 = s71 $\oplus$ s74 \\
s92 = s24 $\oplus$ s61 \\
s93 = s38 $\oplus$ s62 \\
s94 = s58 $\oplus$ s52 \\
s95 = s59 $\oplus$ s63 \\
s96 = s64 $\oplus$ s69 \\
s97 = s61 $\oplus$ s51 \\
s98 = s0 $\oplus$ s65 \\
s99 = s4 $\oplus$ s63 \\
s100 = s58 $\oplus$ s74 \\
s101 = s60 $\oplus$ s50 \\
s102 = s39 $\oplus$ s101 \\
s103 = s34 $\oplus$ s87 \\
s104 = s59 $\oplus$ s89 \\
s105 = s76 $\oplus$ s70 \\
s106 = s92 $\oplus$ s65 \\
s107 = s96 $\oplus$ s98 \\
s108 = s36 $\oplus$ s93 \\
s109 = s95 $\oplus$ s48 \\
s110 = s1 $\oplus$ s92 \\
s111 = s99 $\oplus$ s46 \\
s112 = s98 $\oplus$ s75 \\
s113 = s89 $\oplus$ s90 \\
s114 = s108 $\oplus$ s99 \\
s115 = s110 $\oplus$ s112 \\
s116 = s3 $\oplus$ s104 \\
s117 = s90 $\oplus$ s109 \\
s118 = s108 $\oplus$ s101 \\
s119 = s87 $\oplus$ s112 \\
s120 = s32 $\oplus$ s110 \\
\end{multicols}
\noindent Output signals (output bit $j$ is carried by the listed signal, $j = 0,\dots,31$):
\begin{center}\scriptsize\ttfamily
[96, 119, 83, 113, 118, 78, 79, 84, 107, 103, 76, 116, 102, 80, 85, 91, 106, 115, 105, 117, 114, 81, 86, 94, 97, 120, 77, 109, 111, 82, 88, 100]
\end{center}



\end{document}
